\begin{table}[htbp]
\centering
\caption{Natural Experiment: Colorado's Border with New Mexico}
\label{tab:co_nm}
\small
\begin{tabular}{lcccc}
\toprule
Specification & Coef.\ (NM-border $\times$ post-2022) & SE (HC1) & SE (cl.) & $p$ (cl.) \\
\midrule
Difference-in-differences$^{a}$ & $-0.487$ & $0.106$ & $0.123$ & $<0.001$ \\
\bottomrule
\end{tabular}
\begin{minipage}{0.95\textwidth}
\vspace{0.3em}
\footnotesize
\textit{Notes:} Outcome is log(revenue per active dispensary). Colorado county-year panel, 2014--2025, $n=352$
(2022, the transition year in which New Mexico's market opened partway through, is dropped). Treatment is the
6 Colorado counties within 70 miles of New Mexico (Archuleta, Conejos, Costilla, La Plata, Las Animas,
Montezuma); controls include the 5 Colorado counties bordering a state that remained non-legal throughout the
sample (Utah, Wyoming, Oklahoma, Nebraska, Kansas) and the 31 Colorado interior counties, plus
log(population) and year fixed effects. $^{a}$Post $=1$ for 2023 onward, after New Mexico's April 1, 2022
retail opening. The implied change in revenue per dispensary in New Mexico-border counties after New Mexico's
own market opens is $-38.6$ percent.
\end{minipage}
\end{table}
